\documentclass[11pt]{article}

\usepackage[a4paper,margin=1in]{geometry}
\usepackage{amsmath}
\usepackage{booktabs}
\usepackage{enumitem}
\usepackage{hyperref}

\title{Music Generation with Rule-Based Reinforcement Learning: Literature Notes}
\author{}
\date{\today}

\begin{document}

\maketitle

\section{Project Motivation}
\label{sec:motivation}

The goal of this project is to learn and test whether reinforcement learning,
especially GRPO-style optimization, can improve symbolic music generation when
the reward is based on explicit music-theory rules. The target domain is
classical music, since parts of the style can be formalized more rigidly than in
many other genres.

The core research question is:

\begin{quote}
Can rule-based reinforcement learning improve the structural, harmonic, and
voice-leading quality of symbolic classical music generation without collapsing
diversity or overfitting to narrow heuristics?
\end{quote}

This should be framed as a domain-specific experiment, not as a claim that one
set of music-theory rules can solve general music generation. Rule rewards are
expected to be useful precisely because the initial domain is constrained.

A more concrete target application is Goldberg-style variation generation:

\begin{quote}
Generate a new Bach-like keyboard variation that preserves the Goldberg Aria's
underlying structure while varying the surface melody, rhythm, and texture.
\end{quote}

For this target, ``similarity to the Aria'' cannot mean note-level copying. The
model should preserve deeper structure: meter, 32-bar binary form, phrase
boundaries, harmonic path, bass movement, and cadence locations. At the same
time, the model should be rewarded for producing a genuine variation rather than
reconstructing the Aria.

\section{Papers}
\label{sec:papers}

\subsection{NotaGen: Advancing Musicality in Symbolic Music Generation with Large Language Model Training Paradigms}
\label{sec:paper-notagen}

\textbf{Status: read.}

\textbf{Reference:}
\url{https://arxiv.org/abs/2502.18008}

\textbf{Why it matters.}
NotaGen is the closest modern paper for symbolic classical sheet music
generation. It uses ABC notation, large-scale pretraining, supervised
fine-tuning on high-quality classical compositions, and a preference
optimization stage called CLaMP-DPO.

\textbf{Main points.}
\begin{itemize}[leftmargin=*]
  \item The model is trained on symbolic sheet music, not directly on audio.
  \item Music is represented in ABC notation, which makes symbolic music look
  like structured text.
  \item The paper targets high-quality classical sheet music generation.
  \item The RL/preference stage improves generation quality and controllability,
  but it does not use predefined hand-written music-theory rewards.
  \item It uses bar-stream patching and the Tunesformer architecture.
  \item The most relevant takeaways for this project are the symbolic
  representation, long-piece generation strategy, tokenization design, and
  preference-optimization loop.
\end{itemize}

\textbf{Relevance to this project.}
NotaGen provides a strong modern baseline for symbolic classical generation.
However, our proposed project differs because we want inspectable, rule-based
rewards such as voice-leading, harmony, cadence, range, and meter constraints.

\textbf{Theory links.}
The main technical sections connected to this paper are ABC notation
(Section~\ref{sec:abc}), length/bar annotations (Section~\ref{sec:length}),
bar-stream tokenization (Section~\ref{sec:tokenizers}), and alignment
(Section~\ref{sec:alignment}).

\textbf{Improvement points discussed.}
\begin{itemize}[leftmargin=*]
  \item Replace or supplement CLaMP-style centroid similarity with explicit,
  decomposable symbolic rewards.
  \item Avoid relying on composer-period-instrumentation prompts as if those
  factors were independent.
  \item Treat counterfactual prompts, such as ``Mozart in a Romantic style'', as
  out-of-distribution unless the training setup explicitly supports
  disentanglement.
  \item Track diversity because optimizing toward an average prompt embedding
  can favor central, safe samples.
\end{itemize}

\subsubsection{CLaMP 2 Score and Prompt-Conditioned Similarity}

NotaGen uses the CLaMP 2 Score to quantify semantic similarity between generated
pieces and real pieces associated with the same prompt. Let \(P\) be the set of
prompts. For a prompt \(p \in P\), let \(Y_p\) be the set of real pieces matching
that prompt. CLaMP 2 embeds each real piece into a semantic feature space, and
the average feature vector for the prompt is:

\[
\bar{z}_p = \frac{1}{|Y_p|}\sum_{y \in Y_p} z_y
\]

For a generated piece \(x_p\), the model computes its feature \(z_{x_p}\) and
scores it by similarity to \(\bar{z}_p\). Intuitively, if the prompt is
``Classical period, Mozart, piano'', the generated piece should land near real
pieces matching that prompt in CLaMP 2 embedding space.

This is useful as a high-level style metric, but it raises several concerns:

\begin{itemize}[leftmargin=*]
  \item \textbf{Prompt entanglement.} Composer, period, and instrumentation are
  not independent. For example, the composer often almost determines the period,
  so the model may learn composer identity while treating period as redundant.
  \item \textbf{Weak counterfactual control.} Prompts such as ``Mozart in a
  Romantic style'' may be out-of-distribution if the training set does not
  contain such combinations. The model may ignore one attribute, produce an
  unstable blend, or generate generic music.
  \item \textbf{Centroid-seeking behavior.} Optimizing similarity to the average
  embedding for a prompt can reward safe, central examples. It does not literally
  force the average melody, but strong optimization can reduce novelty and
  diversity by favoring pieces close to the prompt centroid.
\end{itemize}

For this project, CLaMP-style similarity should therefore be treated as a weak
auxiliary signal or evaluation metric, not as the main reward. A cleaner first
objective is:

\[
R = R_{\text{rule compliance}} - \lambda_{\text{KL}}D_{\text{KL}}
  + R_{\text{diversity/novelty}}
\]

This keeps the initial experiment focused on whether explicit symbolic rules can
improve classical generation, while avoiding the harder problem of disentangling
composer, period, and instrumentation.

\subsubsection{CLaMP-DPO Optimization Loop}

The CLaMP-DPO stage is best understood as an iterative generate-score-train
procedure, rather than as online reinforcement learning with an environment. For
each prompt \(p\), the current model generates a candidate set:

\[
X_p = \{x_p^1, x_p^2, \ldots, x_p^n\}
\]

In the paper, \(n\) is approximately 100 pieces per prompt in each iteration.
These are not final outputs; they are a sample pool used to estimate which
generations are better or worse under the CLaMP 2 Score. High-scoring and
low-scoring generations can then be used to construct preference comparisons for
DPO-style training.

The number of outer iterations is denoted by \(K\). NotaGen sets \(K=3\), so the
loop is roughly:

\begin{enumerate}[leftmargin=*]
  \item generate around 100 pieces per prompt with the current model;
  \item embed and score the generated pieces with CLaMP 2;
  \item build preference pairs or rankings from those scores;
  \item update the model with CLaMP-DPO;
  \item repeat for three outer iterations.
\end{enumerate}

The reason for generating many pieces per prompt is that preference optimization
needs comparisons. One sample gives a noisy signal, while a pool of samples gives
enough variation to identify relatively good and bad generations for the same
prompt.

For this project, an analogous loop could be built with rule-based rewards:

\begin{verbatim}
for each prompt or conditioning context:
    generate N candidate pieces
    parse and score each piece with symbolic rules
    form preference pairs from high-score and low-score samples
    train with DPO/GRPO-style optimization
\end{verbatim}

The important design difference is that our score would be decomposable and
inspectable. Instead of rewarding proximity to a style centroid in CLaMP space,
we can directly ask why a sample scored well or poorly: invalid meter, voice
crossing, unresolved dissonance, missing cadence, parallel fifths, low diversity,
and so on.

\subsection{SMART: Tuning a Symbolic Music Generation System with an Audio Domain Aesthetic Reward}
\label{sec:paper-smart}

\textbf{Status: reading.}

\textbf{Reference:}
\url{https://arxiv.org/abs/2504.16839}

\textbf{Why it matters.}
SMART is directly relevant because it uses group relative policy optimization
(GRPO) to fine-tune a symbolic piano MIDI generator.

\textbf{Main points.}
\begin{itemize}[leftmargin=*]
  \item The model generates symbolic piano MIDI.
  \item MIDI is tokenized with the REMI+ tokenizer implemented in MidiTok.
  \item The reward comes from an audio-domain aesthetic model after rendering
  symbolic outputs to audio.
  \item The reward model, Meta Audiobox Aesthetics (MAA), was trained on
  human-rated natural audio and did not use synthetic/generated audio in its own
  training data.
  \item There is no extra alignment fine-tuning stage before GRPO: the policy and
  reference are initialized from the supervised MIDI base model, then GRPO starts
  directly.
  \item GRPO improves average subjective ratings in a small listening study.
  \item The paper reports a key risk: over-optimization can substantially reduce
  diversity.
\end{itemize}

\textbf{Relevance to this project.}
SMART is useful for the RL method and failure modes. The main difference is that
our reward would be explicit and symbolic, rather than an aesthetic model applied
to rendered audio.

\textbf{Theory links.}
The main technical sections connected to this paper are REMI+ tokenization
(Section~\ref{sec:tokenizers}), alignment (Section~\ref{sec:alignment}), and
rule-based rewards (Section~\ref{sec:rewards}).

\textbf{Improvement points discussed.}
\begin{itemize}[leftmargin=*]
  \item MAA is used out-of-distribution: it is trained on natural audio but used
  to score generated MIDI rendered through a soundfont.
  \item The reward is not prompt-conditioned and is not trained on model
  generations, unlike many text RLHF reward/preference models.
  \item The method is a low-effort but useful GRPO template: supervised base
  model, external reward, KL to a frozen reference, and measurement of output
  shifts.
  \item Replace the OOD audio aesthetic reward with on-distribution symbolic rule
  rewards that can be inspected and debugged.
  \item Keep explicit diversity metrics, because SMART shows reward
  over-optimization can collapse output diversity.
\end{itemize}

\subsection{Combining Reinforcement Learning and Constraint Programming for Sequence-Generation Tasks with Hard Constraints}
\label{sec:paper-constraints}

\textbf{Status: read.}

\textbf{Reference:}
\href{https://drops.dagstuhl.de/entities/document/10.4230/LIPIcs.CP.2022.30}{DROPS: 10.4230/LIPIcs.CP.2022.30}

\textbf{Why it matters.}
This is the closest prior work to the rule-based reward direction. It combines
RL with constraint programming for sequence generation, using music generation as
a constrained test domain.

\textbf{Main points.}
\begin{itemize}[leftmargin=*]
  \item Machine learning models imitate a dataset but do not naturally enforce
  hard constraints.
  \item Constraint programming enforces constraints but does not naturally
  imitate a corpus style.
  \item The paper combines both by using constraint information inside an
  RL-based generation setup.
  \item The goal is to generate melodic lines that respect both corpus style and
  music-theory constraints.
  \item The evaluation mainly focuses on constraint satisfaction and corpus
  compatibility, not on broader musical quality.
\end{itemize}

\textbf{Relevance to this project.}
This paper gives a conceptual foundation for using rule or constraint signals.
Our project can modernize the idea using contemporary symbolic models and
GRPO-style training.

\textbf{Theory links.}
The main technical sections connected to this paper are rule-based rewards
(Section~\ref{sec:rewards}), candidate experiments (Section~\ref{sec:experiments}),
and open design choices (Section~\ref{sec:design}).

\textbf{Improvement points to investigate.}
\begin{itemize}[leftmargin=*]
  \item Compare soft scalar rewards with hard constraint satisfaction.
  \item Use modern policy-optimization methods rather than older RL-Tuner style
  methods.
  \item Start with deterministic symbolic rules before adding learned aesthetic
  or style rewards.
  \item Do not copy the evaluation as-is: rule compliance is necessary evidence,
  but it is not enough to show musical quality.
\end{itemize}

\subsection{Music Transformer: Generating Music with Long-Term Structure}
\label{sec:paper-music-transformer}

\textbf{Status: known background.}

\textbf{Reference:}
\url{https://arxiv.org/abs/1809.04281}

\textbf{Why it matters.}
Music Transformer is a foundational symbolic music generation paper. Its main
contribution is not post-training, but the modeling of long-range musical
structure with Transformer attention. It helped establish that symbolic music can
be modeled as an event sequence while still requiring architectural support for
motifs, repetition, rhythmic patterns, and phrase-level dependencies.

\textbf{Main points.}
\begin{itemize}[leftmargin=*]
  \item Uses an event-based symbolic representation rather than raw audio.
  \item Treats music generation as autoregressive sequence modeling.
  \item Emphasizes long-term structure as a central difficulty in music
  generation.
  \item Uses relative attention to better represent relationships between events
  across time.
  \item Shows that local note plausibility is insufficient: the model must
  preserve and develop musical ideas over longer spans.
\end{itemize}

\textbf{Theory links.}
The main technical section connected to this paper is pre-training
(Section~\ref{sec:pretraining}). It also informs the design choices around
tokenization (Section~\ref{sec:tokenizers}) and structure rewards
(Section~\ref{sec:rewards}).

\textbf{Improvement points for this project.}
\begin{itemize}[leftmargin=*]
  \item The base model must be good enough before GRPO; rule-based RL cannot fix
  a model with no musical prior.
  \item Long-range structure should be evaluated separately from local rule
  compliance.
  \item Rule rewards should not only check note-level correctness; they should
  eventually include phrase, repetition, cadence, and development signals.
\end{itemize}

\subsection{Theme Transformer}
\label{sec:paper-theme-transformer}

\textbf{Status: summarized.}

\textbf{Full title:} Theme Transformer: Symbolic Music Generation with
Theme-Conditioned Transformer.

\textbf{Reference:}
\url{https://arxiv.org/abs/2111.04093}

\textbf{Why it matters.}
Theme Transformer is directly relevant to the question of conditioning a music
generator on a source musical idea. The paper argues that prompt-based
conditioning is weak: if a theme is only placed at the beginning of the sequence,
an autoregressive Transformer may gradually ignore it. The authors instead use a
theme-conditioned encoder-decoder architecture, where the theme is encoded
separately and the decoder can cross-attend to it throughout generation.

\textbf{Main idea.}
The paper distinguishes:

\[
p(x_t \mid x_{<t})
\]

for ordinary autoregressive generation, from:

\[
p(x_t \mid x_{<t}; c_{1:\tau})
\]

where \(c_{1:\tau}\) is a separately supplied theme condition. This is important
for our project because the Goldberg Aria skeleton should not merely be a prefix
prompt; it should remain available as an explicit conditioning object during the
whole generation.

\textbf{How themes are obtained.}
The authors do not manually label themes. They automatically retrieve repeated
thematic material from each piece:

\begin{enumerate}[leftmargin=*]
  \item split each song into non-overlapping two-bar fragments;
  \item learn a melody embedding with contrastive learning;
  \item create positive pairs using musically motivated augmentations;
  \item cluster fragments inside a piece with DBSCAN;
  \item use the largest cluster as the theme cluster;
  \item choose the earliest fragment in that cluster as the conditioning theme.
\end{enumerate}

The contrastive augmentations teach the embedding to treat certain surface
changes as still theme-like:

\begin{itemize}[leftmargin=*]
  \item pitch shifting within the scale;
  \item varying the duration of the last note;
  \item splitting or combining repeated notes.
\end{itemize}

\textbf{Architecture.}
The model adds theme-region markers:

\begin{verbatim}
THEME-START
THEME-END
\end{verbatim}

These delimit generated regions where the theme is expected to recur. The
proposed architecture makes two key changes to a standard seq2seq Transformer:

\begin{itemize}[leftmargin=*]
  \item \textbf{Gated parallel attention:} decoder self-attention and
  cross-attention to the theme are arranged in parallel. Inside theme regions,
  upper decoder layers emphasize cross-attention so the theme condition has a
  stronger effect.
  \item \textbf{Theme-aligned positional encoding:} when a theme appears later in
  the generated piece, the positional encoding for cross-attention is reset so
  that position 1 of the generated theme aligns with position 1 of the original
  conditioning theme.
\end{itemize}

\textbf{Data and representation.}
The experiments use POP909, a dataset of Mandarin pop piano arrangements, not
classical Bach. The model uses symbolic MIDI-style event tokens. Notes are
represented with pitch, duration, and velocity tokens, separated for melody and
accompaniment tracks. Metric tokens encode bars, subbeats, and tempo. The
vocabulary contains 730 tokens for the piano representation. Training uses
512-token windows, which correspond to roughly nine bars on average.

\textbf{Evaluation.}
The paper compares a prompt-based Transformer, a basic seq2seq Transformer, and
the proposed Theme Transformer. It evaluates both general coherence and
theme-condition preservation:

\begin{itemize}[leftmargin=*]
  \item pitch-class consistency;
  \item melody inconsistency;
  \item groove consistency;
  \item theme inconsistency;
  \item theme uncontrollability;
  \item theme gap, i.e. the bar distance between theme occurrences.
\end{itemize}

Theme Transformer performs better than the baselines at making the condition
recur with plausible variation. Listening tests also favor it on theme
controllability, theme repetition, and theme variation.

\textbf{Limitations for our project.}
\begin{itemize}[leftmargin=*]
  \item It preserves recurring melodic material, not a harmonic or bass
  scaffold.
  \item The two-bar segmentation is simple and can miss musically meaningful
  phrase boundaries.
  \item The model often creates local variations of the theme rather than a
  convincing long-term development.
  \item The training context is short relative to full-piece form.
  \item The domain is pop piano, not Bach keyboard variation.
\end{itemize}

\textbf{Relevance to the Goldberg direction.}
Theme Transformer supports the core design choice that the Aria structure should
be an explicit condition, not just a textual prompt. However, our conditioning
object should be different from their two-bar melody theme. For Goldberg-style
variation, the condition should be a structured Aria skeleton:

\begin{itemize}[leftmargin=*]
  \item 32-bar binary form;
  \item 3/4 meter;
  \item per-bar bass or scale-degree skeleton;
  \item harmonic progression;
  \item cadence map;
  \item phrase boundaries;
  \item optional variation type or texture target.
\end{itemize}

The project is therefore better described as structure-conditioned variation
generation rather than only theme-conditioned generation. Theme Transformer is a
useful precedent, but our reward should measure preservation of the Aria's deep
structure while penalizing surface copying.

\textbf{Theory links.}
The main technical sections connected to this paper are tokenization
(Section~\ref{sec:tokenizers}), pre-training (Section~\ref{sec:pretraining}),
post-training/alignment (Section~\ref{sec:alignment}), and Goldberg-style
structure rewards (Section~\ref{sec:goldberg-rewards}).

\section{ABC Notation Notes}
\label{sec:abc}

ABC notation is a plain-text representation for symbolic music. It is attractive
for language-model-based music generation because it encodes notes, durations,
voices, meter, key, bar lines, dynamics, and articulations in text.

\subsection{Header and Body}

A minimal ABC example:

\begin{verbatim}
X:1
T:Simple Tune
M:4/4
L:1/4
K:C
C D E F | G A B c |
\end{verbatim}

The header contains metadata:

\begin{itemize}[leftmargin=*]
  \item \verb|X:|: tune or index number.
  \item \verb|T:|: title.
  \item \verb|M:|: meter.
  \item \verb|L:|: default note length.
  \item \verb|K:|: key.
\end{itemize}

The body contains the actual notes and musical events.

\subsection{Durations}

Durations are encoded relative to the default unit length \verb|L:|.

\begin{verbatim}
M:4/4
L:1/8
K:C
C D E2 F/ G3/2 |
\end{verbatim}

If \verb|L:1/8|:

\begin{itemize}[leftmargin=*]
  \item \verb|C| means one eighth note.
  \item \verb|E2| means two units, i.e. a quarter note.
  \item \verb|F/| or \verb|F/2| means half a unit, i.e. a sixteenth note.
  \item \verb|G3/2| means one and a half units, i.e. a dotted eighth note.
\end{itemize}

Rests use the same duration system:

\begin{verbatim}
z      % rest of length L
z2     % rest of length 2 * L
z/2    % rest of length L / 2
\end{verbatim}

\subsection{Bars}

The character \verb+|+ is a bar line. It separates measures:

\begin{verbatim}
M:4/4
L:1/4
K:C
C D E F | G A B c |
\end{verbatim}

In this example, every note has default duration \verb|L:1/4|, so each bar
contains four quarter notes.

Common variants:

\begin{itemize}[leftmargin=*]
  \item \verb+|]+: end of tune.
  \item \verb+||+: double bar line.
  \item \verb+|:+: repeat start.
  \item \verb+:|+: repeat end.
\end{itemize}

\subsection{Voices}

Multiple voices are encoded with \verb|V:| fields and inline voice switches:

\begin{verbatim}
X:1
T:Two voices
M:4/4
L:1/4
K:C
V:1 name="Soprano"
V:2 name="Bass"
[V:1] C D E F | G A B c |
[V:2] C, C, G,, G,, | C, C, G,, G,, |
\end{verbatim}

\verb|[V:1]| switches to voice 1 and \verb|[V:2]| switches to voice 2.
Explicit voices are useful for rule rewards because they allow us to check
voice-leading, voice crossing, ranges, cadences, and parallel intervals.

Chords can be represented with bracketed simultaneous notes:

\begin{verbatim}
[CEG]2
\end{verbatim}

This means C, E, and G played together for duration \verb|2 * L|.

\subsection{Octaves}

Uppercase and lowercase note names indicate different octaves:

\begin{verbatim}
C D E F G A B
c d e f g a b
\end{verbatim}

Lowercase notes are one octave above uppercase notes. Commas lower pitch by
octaves and apostrophes raise pitch:

\begin{verbatim}
C,   % one octave below C
C,,  % two octaves below C
c'   % one octave above c
\end{verbatim}

\subsection{Example Line}

The following line:

\begin{verbatim}
[r:15/30][V:1]._d2 (d2 c2)|
[V:2].f2) (f2 e2)|
[V:3]z2 z2!p! (F,2|
[V:4]z2 z2!p! (F,,2|
\end{verbatim}

can be read as:

\begin{itemize}[leftmargin=*]
  \item \verb|[r:15/30]|: bar-position annotation from the paper, likely current
  bar index 15 and countdown/remaining index 30.
  \item \verb|[V:1]| through \verb|[V:4]|: four voice streams.
  \item \verb|_d2|: D-flat with duration \verb|2 * L|.
  \item \verb|z2|: rest with duration \verb|2 * L|.
  \item \verb|F,2| and \verb|F,,2|: low F notes in different octaves.
  \item \verb|.|: staccato articulation.
  \item \verb|(| and \verb|)|: slur marks.
  \item \verb|!p!|: piano dynamic marking.
\end{itemize}

\section{Length and Bar Annotations}
\label{sec:length}

NotaGen uses stream-based training and inference for long pieces. It annotates
body lines with a marker such as:

\begin{verbatim}
[r:15/30]
\end{verbatim}

The first number is the current bar index. The second is a countdown or remaining
bar index. This is future information about planned length, but not future
information about the actual notes.

The countdown is useful because musical endings depend on knowing that the end
is near. A bar can mean different things depending on the total planned length:

\begin{itemize}[leftmargin=*]
  \item bar 16 of a 16-bar piece is an ending;
  \item bar 16 of a 128-bar piece is early or middle material;
  \item bar 16 of a 32-bar piece is near the midpoint.
\end{itemize}

This makes generation length-conditioned. For a controlled generation setup, the
inference loop can supply these annotations:

\begin{verbatim}
[r:1/31]  generate first bar or line
[r:2/30]  generate second bar or line
...
[r:32/0]  generate ending
\end{verbatim}

For an open-ended generator, countdown annotations are less natural because the
model would need to decide the end itself.

\section{Tokenizers and Symbolic Representations}
\label{sec:tokenizers}

\subsection{Tokenizers Used in the Read Papers}

\subsubsection{NotaGen}

NotaGen uses ABC notation with bar-stream patching. This is a hierarchical
character-level representation:

\begin{itemize}[leftmargin=*]
  \item the score is represented as ABC text;
  \item the text is split into semantically meaningful units, such as header
  lines and bars;
  \item each unit is split into fixed-length character patches and padded when
  needed;
  \item each patch is projected into a patch embedding;
  \item a patch-level decoder models long-range structure;
  \item a character-level decoder predicts the next patch character by
  character.
\end{itemize}

This means the low-level tokens are characters, but the main sequence modeled by
the patch-level Transformer is a sequence of bar/header patch embeddings.

\subsubsection{SMART}

SMART uses MIDI, not ABC. The MIDI files are tokenized using the REMI+ tokenizer
as implemented in the MidiTok library. REMI-style tokenization converts MIDI into
a sequence of musical event tokens, such as:

\begin{verbatim}
Bar
TimeSig_4/4
Tempo_120
Position_0
Pitch_60
Velocity_12
Duration_1.0
\end{verbatim}

The exact vocabulary depends on the MidiTok configuration, but the important
event types are:

\begin{itemize}[leftmargin=*]
  \item \verb|Bar| or similar new-bar tokens;
  \item \verb|TimeSig| tokens for meter;
  \item \verb|Tempo| tokens;
  \item \verb|Position| tokens for onset position inside a bar;
  \item \verb|Pitch| tokens for MIDI pitch;
  \item \verb|Velocity| tokens for loudness;
  \item \verb|Duration| tokens for note length.
\end{itemize}

The paper notes that the tokenizer uses 20 velocity bins. During base model
training, SMART uses random crops up to 512 REMI+ tokens, padded with special
\verb|PAD| tokens when necessary. During GRPO training, prompts are also written
in REMI+ syntax: each prompt includes a token for the start of a new bar, a
random time-signature token, and a random tempo token.

\subsubsection{Theme Transformer}

Theme Transformer uses symbolic MIDI-style event tokens. Each note is represented
by note-related events:

\begin{itemize}[leftmargin=*]
  \item pitch;
  \item duration;
  \item velocity.
\end{itemize}

The representation separates melody and accompaniment tracks, so a melody pitch
and an accompaniment pitch are different token types. It also includes metric
tokens:

\begin{itemize}[leftmargin=*]
  \item bar;
  \item subbeat position inside a bar;
  \item tempo.
\end{itemize}

The distinctive part is the addition of theme-region tokens:

\begin{verbatim}
THEME-START
THEME-END
\end{verbatim}

These markers tell the model where a generated region is intended to instantiate
the conditioning theme. For our Goldberg target, the analogous idea would not be
only theme-region tokens, but explicit structure-conditioning tokens or side
inputs for bar index, Aria bass/harmony, cadence role, and phrase position.

\subsection{BPE}

Byte Pair Encoding (BPE) can be applied to ABC notation. It learns frequent text
chunks, such as:

\begin{verbatim}
" C2"
"|[V:"
"!p!"
"z2 z2"
\end{verbatim}

BPE is easy to implement and efficient in many language-model pipelines. However,
it is not musically aligned by default. It may split or merge musical units in
ways that make symbolic reasoning harder. For rule-based rewards, we still need
to parse the decoded ABC and validate its structure.

\subsection{Bar-Stream Patching}

Bar-stream patching does not mean that a whole bar becomes one vocabulary token.
Instead, it is a hierarchical character-level method.

The pipeline is:

\begin{verbatim}
ABC text
-> split into header lines and bars
-> split each line/bar into fixed-length character patches
-> pad the final patch when needed
-> flatten one-hot character vectors
-> linear projection to one patch embedding
-> patch-level decoder models patch sequence
-> character-level decoder predicts the next patch character by character
\end{verbatim}

For example, a bar:

\begin{verbatim}
[V:1] C2 D2 |
\end{verbatim}

is represented as characters inside a fixed-size patch:

\begin{verbatim}
[ V : 1 ]   C 2   D 2   | <pad> <pad> ...
\end{verbatim}

Each character has an ID or one-hot vector. The full fixed-length character
array is projected into one patch embedding.

\subsection{Training Objective}

The model predicts characters, not patch IDs. For target patch \(P_{i+1}\):

\[
p(P_{i+1} \mid P_1, \ldots, P_i)
= \prod_t p(c_t \mid P_1, \ldots, P_i, c_{<t})
\]

During training, teacher forcing provides the previous characters in the current
target patch. A causal mask prevents each character position from seeing future
characters, so all character losses can still be computed in parallel.

The context for generating a character inside the current patch is:

\begin{verbatim}
compressed previous patches + previous characters within current patch
\end{verbatim}

It is not:

\begin{verbatim}
all previous characters from the full piece as separate tokens
\end{verbatim}

This is the key architectural tradeoff. Bar-stream patching improves efficiency
and gives the model a bar-level structural bias, but earlier patch internals are
compressed into patch embeddings.

\section{Pre-Training}
\label{sec:pretraining}

The project should separate two questions:

\begin{enumerate}[leftmargin=*]
  \item Can we train or reuse a symbolic model that already has a useful musical
  prior?
  \item Can post-training with rule-based rewards improve that model without
  damaging style or diversity?
\end{enumerate}

Music Transformer is the relevant foundational background for the first
question. Its lesson is that symbolic music generation is not only next-token
prediction. A model can produce locally plausible notes while failing at
long-range musical coherence.

\subsection{Autoregressive Symbolic Modeling}

The pre-training objective is standard language modeling over symbolic music
tokens:

\[
p(x) = \prod_t p(x_t \mid x_{<t})
\]

The tokens may be ABC characters, ABC patches, REMI+ MIDI events, or more
music-aware note/duration/chord events. In all cases, the base model learns a
distribution over musical sequences from data before any rule-based RL is
applied.

For this project, pre-training or supervised fine-tuning should provide:

\begin{itemize}[leftmargin=*]
  \item basic syntax validity;
  \item plausible local note transitions;
  \item common rhythmic patterns;
  \item style priors from the target corpus;
  \item some ability to maintain meter and phrase-level expectations.
\end{itemize}

Rule-based GRPO should be viewed as post-training on top of this prior, not as a
replacement for learning musical style from data.

\subsection{Long-Term Structure}

Music Transformer emphasizes that musical quality depends on relationships over
long spans:

\begin{itemize}[leftmargin=*]
  \item motif repetition and variation;
  \item rhythmic consistency;
  \item phrase structure;
  \item harmonic direction;
  \item preparation and resolution;
  \item cadence placement.
\end{itemize}

This is important for our reward design. A reward that only checks local
constraints can make music more formally correct while leaving it boring or
fragmented. Long-range structure should therefore be evaluated separately from
local rule compliance.

\subsection{Representation Implications}

The representation chosen for pre-training affects what the model can learn and
what the reward can inspect:

\begin{itemize}[leftmargin=*]
  \item ABC text is convenient for sheet music and classical symbolic structure.
  \item Bar-stream patching gives the model a structural bias around bars while
  retaining character-level generation.
  \item REMI+ MIDI event tokens are convenient for performance-like piano MIDI,
  especially when velocity, tempo, and duration matter.
  \item Music-aware symbolic tokens may be simpler for a first Bach/SATB reward
  pipeline because they map directly to notes, durations, voices, and bars.
\end{itemize}

The first implementation should prefer a representation that is easy to parse and
score. A slightly less efficient tokenizer is acceptable if it makes the reward
module reliable.

\subsection{Base Model Requirement}

The base model must be strong enough that generated candidates contain meaningful
musical variation. GRPO can then prefer better candidates within a group. If all
candidate samples are invalid or musically empty, the reward signal will mostly
teach syntax hacks rather than musical improvement.

A practical pre-training baseline is:

\begin{verbatim}
target corpus: Bach chorales or another constrained classical subset
representation: ABC, MusicXML-derived events, or REMI-like symbolic events
objective: next-token prediction
evaluation: syntax validity, bar validity, corpus likelihood, diversity
\end{verbatim}

Only after this baseline generates mostly parseable and stylistically plausible
samples should we move to post-training with rule rewards.

\section{Post-Training and Alignment}
\label{sec:alignment}

The two alignment methods discussed so far are quite different. NotaGen uses an
offline preference-optimization loop based on CLaMP-DPO, while SMART uses direct
on-policy GRPO with an external audio aesthetic reward. The former is more
interesting as an alignment pipeline, even though its reward target is not ideal
for our project. The latter is a useful minimal implementation template for GRPO.

\subsection{NotaGen: CLaMP-DPO}

NotaGen starts from a pretrained and supervised fine-tuned symbolic model. It
then uses CLaMP 2 as an automatic scorer to construct preference data.

The method is:

\begin{enumerate}[leftmargin=*]
  \item define prompts \(p\), such as period-composer-instrumentation
  combinations;
  \item compute a reference embedding centroid \(\bar{z}_p\) from real pieces
  matching each prompt;
  \item generate a candidate set \(X_p\) for each prompt;
  \item score generated pieces by similarity between \(z_{x_p}\) and
  \(\bar{z}_p\);
  \item construct preference pairs from relatively high-scoring and low-scoring
  generations;
  \item update the model with DPO-style training.
\end{enumerate}

The paper sets \(K=3\) outer iterations and generates approximately 100 pieces
per prompt in each iteration. The 100 samples are used to create enough
within-prompt variation for preference construction; they are not the final
output count.

\textbf{Strengths.}
\begin{itemize}[leftmargin=*]
  \item Uses generated samples to construct alignment data.
  \item The comparison is prompt-specific rather than global.
  \item Avoids training a new human preference model.
  \item Fits well with symbolic generation because CLaMP operates on symbolic
  music representations.
\end{itemize}

\textbf{Weaknesses for our project.}
\begin{itemize}[leftmargin=*]
  \item The score is similarity to an embedding centroid, not explicit musical
  correctness.
  \item Prompt variables are entangled; composer may dominate period.
  \item Counterfactual prompts are not naturally supported.
  \item Strong optimization can favor typical, central samples over creative or
  unusual but valid music.
\end{itemize}

\textbf{Adaptation idea.}
Use the same generate-score-preference-train structure, but replace CLaMP
similarity with decomposable symbolic rewards:

\begin{verbatim}
generate candidates
parse symbolic music
score meter, harmony, voice leading, cadence, diversity
construct high-vs-low preference pairs
train with DPO or GRPO-style optimization
\end{verbatim}

\subsection{SMART: Direct GRPO with MAA}

SMART starts from a supervised MIDI base model. The policy and frozen reference
model are initialized from that same base model, and then the system directly
runs GRPO.

The method is:

\begin{enumerate}[leftmargin=*]
  \item generate REMI+ prompts procedurally, including bar, time-signature, and
  tempo tokens;
  \item sample MIDI completions from the policy model;
  \item render each MIDI completion to audio using a soundfont renderer;
  \item crop audio to 10 seconds, matching MAA's receptive field;
  \item score each audio sample with Meta Audiobox Aesthetics;
  \item compute group-relative advantages from the rewards;
  \item update the policy with a KL penalty to the frozen reference model.
\end{enumerate}

This can be summarized as:

\begin{verbatim}
supervised MIDI LM
-> policy/reference initialization
-> generate MIDI
-> render to audio
-> score with MAA
-> GRPO update with KL to reference
\end{verbatim}

\textbf{Strengths.}
\begin{itemize}[leftmargin=*]
  \item Simple GRPO recipe.
  \item Requires no custom reward-model training.
  \item Demonstrates that RL can measurably shift symbolic music outputs.
  \item Provides clear warnings about diversity collapse under stronger
  optimization.
\end{itemize}

\textbf{Weaknesses for our project.}
\begin{itemize}[leftmargin=*]
  \item MAA was trained on natural human-rated audio, not on generated/rendered
  symbolic music.
  \item The reward model is not prompt-conditioned.
  \item The reward is only 10-second audio-level aesthetics, so it cannot capture
  long-range musical form.
  \item The reward can be hacked through low-level correlates such as note
  density, fewer pauses, more dynamics, or wider pitch range.
  \item There is no intermediate generated-output preference dataset or
  alignment fine-tuning stage.
\end{itemize}

\textbf{Adaptation idea.}
Keep the simple GRPO skeleton, but score the generated symbolic object directly:

\begin{verbatim}
supervised symbolic LM
-> policy/reference initialization
-> generate symbolic candidates
-> parse symbolic candidates
-> score deterministic music-theory rules
-> GRPO update with KL to reference
\end{verbatim}

This removes the audio rendering and OOD aesthetic reward from the critical path.
The reward becomes narrower, but it is inspectable and on-distribution.

\subsection{Constraint Paper: Constraints as Modern GRPO Rewards}

The constraint-programming paper is useful less as an implementation template and
more as a source of constraint ideas. We do not need to inherit the full CP
machinery in the first version of this project.

The practical takeaway is:

\begin{verbatim}
music generation can be guided by explicit constraints
constraints can be converted into reward signals
constraint satisfaction alone is not enough as evaluation
\end{verbatim}

A modern post-training version can be much simpler:

\begin{verbatim}
base symbolic model
-> generate multiple candidates per prompt/context
-> score each candidate with explicit symbolic rule functions
-> use GRPO to increase probability of higher-scoring candidates
-> keep KL to a reference model to preserve style/diversity
\end{verbatim}

In this setup, the reward can be implemented with ordinary symbolic checks rather
than a CP solver:

\[
\begin{aligned}
R ={}& R_{\text{parse}}
+ R_{\text{bar duration}}
+ R_{\text{voice range}} \\
&- P_{\text{voice crossing}}
- P_{\text{parallel fifths/octaves}}
- P_{\text{unresolved dissonance}} \\
&+ R_{\text{cadence}}
- P_{\text{repetition}}
\end{aligned}
\]

This gives us a cleaner first implementation:

\begin{itemize}[leftmargin=*]
  \item simpler than integrating a CP solver;
  \item directly compatible with GRPO-style post-training;
  \item interpretable, because each reward component can be inspected;
  \item easier to debug when the model starts reward hacking;
  \item easy to extend toward Bach chorale or SATB-specific rules.
\end{itemize}

CP may become useful later if we need hard guarantees, constrained decoding,
repair of invalid samples, or global feasibility estimates for partial
generations. For the first milestone, plain rule functions are the better
engineering choice.

\section{Rule-Based Reward Ideas}
\label{sec:rewards}

The reward can be built from symbolic checks over parsed ABC or converted
MusicXML/MIDI representations. A possible reward decomposition is:

\[
R = R_{\text{parse}} + R_{\text{meter}} + R_{\text{harmony}}
  + R_{\text{voice}} + R_{\text{cadence}} + R_{\text{style}}
  - P_{\text{degenerate}}
\]

\subsection{Goldberg-Style Aria Structure Rewards}
\label{sec:goldberg-rewards}

For Goldberg-style variation generation, the main reward should not be generic
similarity to Bach or raw similarity to the Aria. The objective should preserve
the Aria's underlying structure while encouraging a different surface.

A useful conditioning object is:

\begin{itemize}[leftmargin=*]
  \item 32-bar binary form;
  \item 3/4 meter;
  \item per-bar bass or scale-degree skeleton;
  \item harmonic progression or chord labels;
  \item cadence map;
  \item phrase boundaries;
  \item optional variation type, such as canon, figuration, dance-like texture,
  or toccata-like texture.
\end{itemize}

A first reward decomposition could be:

\[
\begin{aligned}
R ={}& R_{\text{parse}}
+ R_{\text{meter}}
+ R_{\text{bar count}}
+ R_{\text{cadence map}} \\
&+ R_{\text{bass/harmony similarity}}
+ R_{\text{keyboard playability}}
+ R_{\text{Bach style}} \\
&- P_{\text{surface copy}}
- P_{\text{degenerate repetition}}
- P_{\text{invalid texture}}
\end{aligned}
\]

The important distinction is:

\begin{itemize}[leftmargin=*]
  \item reward similarity in harmony, bass function, phrase location, and cadence
  placement;
  \item penalize direct note-level copying of the Aria;
  \item leave room for different rhythmic figuration, contrapuntal texture, and
  melodic surface.
\end{itemize}

For DPO, candidate outputs for the same Aria-conditioned prompt can be ranked by
this structure-aware reward. For GRPO, the same reward can be applied directly
to a group of generated candidates. This is the concrete version of the project:
post-train a symbolic model to preserve deep musical structure while varying the
surface.

\subsection{Low-Level Validity}

\begin{itemize}[leftmargin=*]
  \item Reward parseable ABC.
  \item Penalize malformed headers, unmatched slurs, invalid voice markers, and
  invalid duration syntax.
  \item Reward every bar having the correct total duration under \verb|M:| and
  \verb|L:|.
  \item Penalize missing or extra bar lines.
\end{itemize}

\subsection{Voice-Level Rules}

\begin{itemize}[leftmargin=*]
  \item Keep voices within plausible SATB ranges.
  \item Penalize voice crossing.
  \item Penalize excessive leaps.
  \item Reward stepwise melodic motion where stylistically appropriate.
  \item Penalize unresolved leading tones.
\end{itemize}

\subsection{Harmony and Counterpoint}

\begin{itemize}[leftmargin=*]
  \item Reward consonant intervals on strong beats.
  \item Penalize parallel fifths and octaves.
  \item Penalize unresolved dissonances.
  \item Reward plausible cadences near phrase endings.
  \item Reward final tonic closure in simple classical settings.
\end{itemize}

\subsection{Structure}

\begin{itemize}[leftmargin=*]
  \item Use target-length conditioning or countdown annotations to encourage
  phrase endings.
  \item Reward phrase-level regularity, such as 4-bar or 8-bar phrases.
  \item Penalize repeated trivial loops and excessive copying.
\end{itemize}

\section{Candidate Experiments}
\label{sec:experiments}

\subsection{Experiment 1: Rule Compliance After Supervised Fine-Tuning}

Train a small symbolic model on a constrained dataset, such as Bach chorales.
Evaluate whether the model already satisfies basic rule metrics after
supervised training.

\textbf{Goal:} establish baseline rule violations before RL.

\subsection{Experiment 2: GRPO With Parse and Meter Rewards}

Start with low-risk rewards:

\begin{itemize}[leftmargin=*]
  \item parseable ABC;
  \item correct bar duration;
  \item valid voice markers;
  \item no empty or degenerate bars.
\end{itemize}

\textbf{Goal:} test the RL pipeline before using subjective or complex music
rules.

\subsection{Experiment 3: GRPO With Voice-Leading Rewards}

Add SATB-specific rewards:

\begin{itemize}[leftmargin=*]
  \item range constraints;
  \item voice crossing penalties;
  \item parallel fifth and octave penalties;
  \item melodic leap penalties.
\end{itemize}

\textbf{Goal:} test whether explicit classical rules improve symbolic quality.

\subsection{Experiment 4: Reward Hacking and Diversity}

Measure diversity before and after RL:

\begin{itemize}[leftmargin=*]
  \item pitch-class entropy;
  \item n-gram diversity over notes or bars;
  \item unique rhythmic patterns;
  \item duplicate output rate;
  \item average edit distance between generations.
\end{itemize}

\textbf{Goal:} avoid the over-optimization problem reported in SMART.

\subsection{Experiment 5: Tokenization Comparison}

Compare simple alternatives:

\begin{itemize}[leftmargin=*]
  \item flat character-level Transformer;
  \item music-aware tokenization;
  \item BPE;
  \item bar-stream patching if implementation time allows.
\end{itemize}

\textbf{Goal:} understand whether bar-aligned structure matters for rule-based
optimization.

\section{Open Design Choices}
\label{sec:design}

\begin{itemize}[leftmargin=*]
  \item Dataset: Bach chorales are likely the cleanest starting point because
  SATB rules are well defined.
  \item Representation: ABC is useful for language modeling; MusicXML or
  music21 objects may be better for reward computation.
  \item Tokenization: character-level or music-aware tokens are simpler than
  implementing full bar-stream patching at first.
  \item Conditioning: use target length or countdown annotations if we want
  planned phrases and cadences; avoid them if we want open-ended generation.
  \item Reward design: start with validity and meter before adding complex
  harmony and counterpoint.
  \item Evaluation: combine objective rule metrics with listening or expert
  preference checks later.
\end{itemize}

\section{Initial Project Direction}
\label{sec:initial-direction}

A practical first milestone is:

\begin{enumerate}[leftmargin=*]
  \item Build or load a small Bach chorale dataset.
  \item Convert pieces into a stable symbolic representation.
  \item Implement a parser and rule-evaluation module.
  \item Train a small supervised baseline.
  \item Verify the baseline has a useful musical prior: parseable syntax,
  plausible local motion, and nontrivial diversity.
  \item Run GRPO or a simpler post-training objective with parse and meter
  rewards.
  \item Add voice-leading rewards only after the pipeline is stable.
\end{enumerate}

This keeps the project focused: prove that rule-based optimization can improve
measurable symbolic correctness before trying to optimize broad musical quality.

\end{document}
